 {\bfseries Important note on how to use this manual}

This short user manual is meant to provide the user with enough information to understand the logic and various options of Ge\+N-\/\+Foam. However, consistent with the Doxygen philosophy, and with the objective of minimizing inconsisitencies between documentation and cource code, detailed usage information of non-\/trivial sub-\/solvers, behavioural models, etc., are (or will be) included directly in the header (.H) file of the corresponsing classes. This manual provides links to most of these header files. As an alternative, one can search them by using the search function at the top right of the page. The links provided in the documentation will bring you to the Doxygen page of the corresponding header file. In this page, one can find the mentioned usage information under the section {\itshape Classes}. It may happen that the header file was not appropriately formatted for Doxygen at the time of its creation, in which case one may have to directly look at the .H file to find the usage information (work in progress to avoid that). ~\newline
~\newline
 In a similar fashion, to help describe the use of complex dictionaries (i.\+e., input files), this manual provides links to one or more commented dictionaries that are available in the tutorials. ~\newline
~\newline
 Exact keywords for sub-\/solvers, models, etc can be found in the corresponding header files. However an easier method to find this names consists in the classical Open\+F\+O\+AM Banana method\+: write in the dictionary \char`\"{}\+Banana\char`\"{} (or any funny word you like), and Ge\+N-\/\+Foam will normally give you an error and a list of (typically self-\/explanatory) valid keywords. 

 N.\+B.\+: {\bfseries Users are expected to be already familiar with Open\+F\+O\+AM and nuclear engineering!}. 

\subsection*{Some theoretical background}

A fairly general theoretical presentation of Ge\+N-\/\+Foam is provided in Ref. \cite{FIORINA201524}. It is recommended to go through this paper before starting to use Ge\+N-\/\+Foam. However, the paper is getting quite old and it is recommended to refer to Ref. \cite{FIORINA2016212} for the diffusion solver, Ref. \cite{FIORINA2017419} for the S\+P3 solver, Refs. \cite{Fiorina2019DetailedOpenFoam} \cite{Fiorina2015ApplicationCodes} for the thermal-\/mechanic solver and its use for mesh deformation, Ref. \cite{Fiorina2019DetailedOpenFoam} for the SN solver, and Refs. \cite{Radman2019ADesign} \cite{RADMAN2021111178} \cite{RADMAN2021111422} for single-\/ and two-\/phase thermal-\/hydraulics.

Before using this manual, we reccomend to go through the introductury lectures to both Open\+F\+O\+AM and Ge\+N-\/\+Foam that are provided in the folder {\itshape Documentation/some\+Useful\+Documents\+And\+P\+Resentations}. These lectures are taken from an I\+A\+EA e-\/learning course available at \href{https://elearning.iaea.org/m2/course/view.php?id=1286}{\tt https\+://elearning.\+iaea.\+org/m2/course/view.\+php?id=1286}. The course requires registration and a N\+U\+C\+L\+E\+US account, but it should be available to all I\+A\+EA member states.

\subsection*{Some practical information}

Here below a couple of essential points that make Ge\+N-\/\+Foam different than most of the other Open\+F\+O\+A\+M-\/based solvers.

{\bfseries The multi-\/region approach}

Ge\+N-\/\+Foam employs a multi-\/region approach to model different physics using different meshes. This implies that the {\itshape 0}, {\itshape constant} and {\itshape system} folders of each case contains multiple folders, one for each physics. In particular, the regions {\itshape fluid\+Region}, {\itshape neutro\+Region} and {\itshape thermo\+Mechanical\+Region} are employed in Ge\+N-\/\+Foam for thermal-\/hydraulics, neutronics and thermal-\/mechanics. There is no requirement for the three meshes to occupy the same region of space. Consistent mapping of fields is performed and a reference value is given to a field if no correspondence is found in the mesh where its value is being projected from.

 {\bfseries The meshes}

A dummy mesh must always be present in all physics (region) directories, even if not solved for. The E\+M\+P\+TY case is already provided with minimal dummy meshes and consistent fields in the “0” folder. Be careful! In case of parallel calculations all your meshes will have to have a number of cells equal or higher than the number of domains you are decomposing your geometry into. In case you need more cells than what available in the E\+M\+P\+TY case, you can run a refine\+Mesh. 

{\bfseries The multi-\/zone approach}

In order to assign different properties (for instance, different porous medium properties or different cross sections) to different zones in a mesh, Ge\+N-\/\+Foam employs the Open\+F\+O\+AM concept of cell\+Zone. Each mesh should then be divided in different cell\+Zones. Each cell\+Zone is associated to a name and this name is used in {\itshape constant/fluid\+Region/phase\+Properties}, {\itshape constant/neutro\+Region/nuclear\+Data}, {\itshape constant/themo\+Mechanical\+Region/themo\+Mechanical\+Properties} to associate each cell\+Zone with a set of properties. The creation of cell\+Zones is normally allowed by all meshers, though different names are normally used (for instance, {\itshape physical entities} in gmsh and {\itshape groups} in Salome). In some cases, conversion of the mesh into an Open\+F\+O\+AM format creates cell\+Set instead of cell\+Zones. In these case, one can use the topo\+Set utility to convert cell\+Sets into cell\+Zones.

\subsection*{The source code}

Ge\+N-\/\+Foam is an open-\/source code and makes use in its programming of fairly high-\/level A\+PI, intuitive naming of variables, and frequent comments. As such, we encourage users to consider the code itlsef as an essential part of the documentation.

The source code is subdivided into 3 main folders\+:
\begin{DoxyItemize}
\item main\+: containing the main \hyperlink{GeN-Foam_8C}{Ge\+N-\/\+Foam.\+C} source file (and other files directly employed by it), which is nothing but a fairly complex coupling loop that calls various functionalities that are found under {\itshape neutronics}, {\itshape thermal\+Hydraulics}, and {\itshape thermo\+Mechanics};
\item classes\+: containing the 3 main classes (or sub-\/libraries of classes) employed to solve for neutronics, thermal-\/hydraulics and thermal-\/mechanics, as well as a class for multi-\/physics controls;
\item include\+: containing specialized versions of some Open\+F\+O\+AM base functionalities.
\end{DoxyItemize}

As a general rule, most classes have an {\itshape include} folder that is used to store all the .H files that are included via \char`\"{}\#include\char`\"{} in the definition of the respective class (normally in the .C file). This is done only to avoid extremely long .C files.

\subsection*{Content of this manual}

The following sections describe the use of the 3 main sub-\/solvers of Ge\+N-\/\+Foam, incuding some theory and useful references\+:
\begin{DoxyItemize}
\item \hyperlink{NEUTRONICS}{Neutronics}
\item \hyperlink{TH}{Thermal-\/hydraulics}
\item \hyperlink{TM}{Thermal-\/mechanics}
\end{DoxyItemize}

The following section describes instead the coupling strategy and the general Ge\+N-\/\+Foam options.


\begin{DoxyItemize}
\item \hyperlink{COUPLING}{Coupling options and time stepping}
\end{DoxyItemize}

© All rights reserved. E\+C\+O\+LE P\+O\+L\+Y\+T\+E\+C\+H\+N\+I\+Q\+UE F\+E\+D\+E\+R\+A\+LE DE L\+A\+U\+S\+A\+N\+NE, Switzerland, 2021 